% Adición focal propuesta; no sustituye el álgebra electrónica existente.
% Inserción: después de la proposición «Álgebra del bloque electrónico local»
% y su párrafo explicativo, antes de «Transferencia entre hojas...».
% Procedencia y alcance: LOCALIZADORES_Y_ALCANCE.md, en esta carpeta.
\subsection{Representación espinorial del bloque electrónico y helicidad}
\label{sec:el-representacion-espinorial}

La estructura electrónica dispone de una representación que conserva más
información que su evaluación escalar. El plano principal de los dos
proyectores de APP tridimensional, seleccionado por el modo trítico
\((1,-1,0)^3\), ya ha proporcionado los operadores \(S\) y \(J\).
El primero es una involución simétrica; el segundo es el conmutador
normalizado y conserva la orientación de las hojas. Su composición permite
construir explícitamente una representación de rotaciones y separar sus
dos componentes de helicidad. Esta construcción actúa sobre el bloque
electrónico; el registro TPK de la ruta y su memoria permanecen como datos
de procedencia del bloque.

Sea \(E_e\) el plano real principal, dotado del producto escalar heredado
de las esperanzas condicionales, y sea
\(\mathscr S_e=E_e\otimes_{\mathbb R}\mathbb C\) su complejificación.
Los operadores ya obtenidos satisfacen
\[
 S^*=S,\qquad J^*=-J,\qquad
 S^2=I,\qquad J^2=-I,\qquad SJ=-JS.
\]
La complejificación extiende el cuerpo de representación. Conserva los
operadores reales y permite distinguir la unidad escalar \(i\) del
endomorfismo \(J\), aunque ambos tengan cuadrado \(-1\) en sus
respectivos dominios.

\begin{proposition}[Terna hermítica inducida por las dos hojas]
\label{prop:el-terna-hermitica}
Los operadores
\begin{equation}
 \sigma_1=SJ,\qquad \sigma_2=-iJ,\qquad \sigma_3=S
 \label{eq:el-terna-hermitica}
\end{equation}
son hermíticos y de traza nula. Satisfacen
\begin{equation}
 \sigma_a\sigma_b
 =\delta_{ab}I+i\sum_{c=1}^3\epsilon_{abc}\sigma_c,
 \qquad
 \frac12\operatorname{tr}(\sigma_a\sigma_b)=\delta_{ab},
 \label{eq:el-algebra-terna}
\end{equation}
donde \(\epsilon_{123}=1\). Por tanto constituyen una base ortonormal
del espacio real \(V_e\) de endomorfismos hermíticos de traza nula de
\(\mathscr S_e\), para el producto
\(\langle X,Y\rangle=\frac12\operatorname{tr}(XY)\).
\end{proposition}
\begin{proof}
Las relaciones anteriores dan \((SJ)^*=(-J)S=SJ\) y
\((-iJ)^*=-iJ\). Los cuadrados de las tres matrices son \(I\).
Además,
\[
 (SJ)(-iJ)=iS,\qquad (-iJ)S=iSJ,\qquad S(SJ)=J=i(-iJ).
\]
Al invertir el orden cambia el signo. Esto prueba la primera identidad
de \eqref{eq:el-algebra-terna}. En la base ortonormal del plano principal,
las matrices son
\[
 \sigma_1=\begin{pmatrix}0&1\\1&0\end{pmatrix},\quad
 \sigma_2=\begin{pmatrix}0&-i\\i&0\end{pmatrix},\quad
 \sigma_3=\begin{pmatrix}1&0\\0&-1\end{pmatrix}.
\]
Su traza es nula, y tomar trazas en la identidad de producto prueba la
ortonormalidad. Todo endomorfismo hermítico de traza nula tiene la forma
\(\bigl(\begin{smallmatrix}z&x-iy\\x+iy&-z\end{smallmatrix}\bigr)\),
con \(x,y,z\) reales, por lo que pertenece a su envolvente real.
\end{proof}

Esta terna recibe, en la representación matricial usual, el nombre de
matrices de Pauli. Su empleo aquí es posterior a la construcción de
\(P,Q,S,J\): se obtiene por los productos
\eqref{eq:el-terna-hermitica}. El espacio de direcciones de esta
representación es la esfera unidad de \(V_e\), cuya forma positiva
acaba de construirse mediante la traza.

\begin{proposition}[Generadores de rotación y descomposición direccional]
\label{prop:el-helicidad}
Defínanse \(L_a=\sigma_a/2\). Entonces
\begin{equation}
 [L_a,L_b]=i\sum_c\epsilon_{abc}L_c,
 \qquad \sum_{a=1}^3L_a^2=\frac34 I.
 \label{eq:el-generadores-spin}
\end{equation}
Para \(n=(n_1,n_2,n_3)\) con \(\sum_an_a^2=1\), sean
\[
 H_e(n)=\sum_an_a\sigma_a,\qquad
 h_e(n)=\frac12H_e(n),\qquad
 \Pi_\pm(n)=\frac12\bigl(I\pm H_e(n)\bigr).
\]
La fórmula lineal \(H_e(v)=\sum_av_a\sigma_a\) se utiliza también
para \(v\) de norma arbitraria.
Se tiene
\begin{equation}
 \mathscr S_e=\operatorname{im}\Pi_+(n)
             \mathbin{\oplus}\operatorname{im}\Pi_-(n),
 \qquad
 h_e(n)\Pi_\pm(n)=\pm\frac12\Pi_\pm(n).
 \label{eq:el-descomposicion-helicidad}
\end{equation}
Los dos sumandos son líneas complejas ortogonales. El transporte de
rotación alrededor de \(n\) es
\begin{equation}
 U_n(\theta)=\exp\!\left(-\frac{i\theta}{2}H_e(n)\right)
 =\cos\frac\theta2\,I-i\sin\frac\theta2\,H_e(n),
 \qquad
 U_n(2\pi)=-I,\quad U_n(4\pi)=I.
 \label{eq:el-retorno-spin}
\end{equation}
\end{proposition}
\begin{proof}
La identidad de producto de la terna da los conmutadores de
\eqref{eq:el-generadores-spin}; sus cuadrados suman \(3I/4\).
Los términos antisimétricos se cancelan al multiplicar
\(H_e(n)\) por sí mismo, de modo que \(H_e(n)^2=I\).
Se deduce
\[
 \Pi_\pm(n)^2=\Pi_\pm(n),\qquad
 \Pi_+(n)\Pi_-(n)=0,\qquad \Pi_+(n)+\Pi_-(n)=I.
\]
Son hermíticos y tienen traza \(1\); por ello su rango es \(1\).
Multiplicar por \(H_e(n)/2\) prueba
\eqref{eq:el-descomposicion-helicidad}. La serie de la exponencial se
separa en potencias pares e impares gracias a \(H_e(n)^2=I\), lo que
produce \eqref{eq:el-retorno-spin}. También prueba directamente la
unitariedad y la ley \(U_n(\theta+\psi)=U_n(\theta)U_n(\psi)\).

La normalización angular se verifica sobre \(V_e\). Para todo
\(v\in\mathbb R^3\), la misma identidad de producto da
\[
 U_n(\theta)H_e(v)U_n(\theta)^*
 =H_e\!\left(v\cos\theta+(n\mathbin{\times}v)\sin\theta
       +n(n\mathbin{\cdot}v)(1-\cos\theta)\right).
\]
La acción inducida es una rotación de ángulo \(\theta\) y eje \(n\).
El factor \(1/2\) en el generador espinorial es, por tanto, el que
corresponde a esa parametrización angular. Una rotación completa restaura
la dirección en \(V_e\) y cambia el signo del vector de
\(\mathscr S_e\); dos rotaciones completas restauran ambos.
\end{proof}

La representación es irreducible: un subespacio complejo invariante por
las tres matrices sería invariante por toda \(M_2(\mathbb C)\), el
álgebra que generan. Sus únicos subespacios invariantes son \(0\) y
\(\mathscr S_e\). Las relaciones
\eqref{eq:el-generadores-spin} identifican así la representación de
espín \(1/2\). Si \(n\) es la dirección de propagación en una
realización cinemática, \(h_e(n)\) es su helicidad adimensional, con
valores \(\pm1/2\). La dirección cinemática se declara al efectuar esa
realización; los dos proyectores están definidos para cada dirección
antes de seleccionar un estado de propagación particular.

\paragraph{Conservación de la estructura electrónica y tipos de inversión.}
La inversión de dirección tiene una ley exacta:
\[
 H_e(-n)=-H_e(n),\qquad \Pi_\pm(-n)=\Pi_\mp(n).
\]
Esta operación intercambia las etiquetas de helicidad relativas a una
dirección. La conjugación de la carga actúa, en cambio, sobre la firma
orientada de la ruta material; la inversión celular actúa sobre las
coordenadas del atlas. Sus dominios y acciones permanecen distintos.
La unicidad de la celda \((5,5)\) es compatible con las dos líneas de
\eqref{eq:el-descomposicion-helicidad}: una celda del atlas no es un
vector de la representación espinorial que porta.

La evaluación electrónica \(\varepsilon_e^\beta\) conserva su
fórmula y sus lectores. Sobre \(\mathscr S_e\), su realización escalar
es \(\varepsilon_e^\beta I\), que conmuta con los generadores y con
\(\Pi_\pm(n)\). La descomposición de helicidad añade la información
rotacional sin modificar la evaluación ni convertir el espín en un
coeficiente multiplicativo de la masa.

La quiralidad corresponde a la graduación de la representación relativista
que se utilice: una involución \(\Gamma_{\rm ch}\) determina sus
proyectores \((I\mp\Gamma_{\rm ch})/2\). La helicidad de
\eqref{eq:el-descomposicion-helicidad} depende además de \(n\).
Por su parte, la holonomía de una banda real de Möbius describe el cambio
de signo de una línea al recorrer un lazo de su base. La ley de esa línea,
el retorno espinorial \eqref{eq:el-retorno-spin} y la holonomía nonádica
de fase con avance de memoria son leyes de transporte sobre objetos
especificados. Su comparación conserva esos objetos: la forma helicoidal
de una trayectoria no sustituye la representación que acaba de
construirse, ni fija por sí misma la quiralidad o la carga.
